\subsubsection{Primitives and Ordinary Differential Equations}\label{subsec:distribution-primitives-odes}

If the derivatives of a function vanish, the function is a constant. We shall now establish the analogous result for distributions.

\begin{Theorem}
\label{thm:distribution-zero-gradient-constant}
Let $\Omega$ be connected, and let $u\in\mathcal{D}^{\prime}(\Omega)$ be such that $\nabla u=0$. Then $u$ is a constant.
\end{Theorem}

\begin{Proof}
We first consider the one-dimensional case. Let $\Omega=I$ be an interval. The condition that $u'=0$ means that $(u,\phi')=0$ for every $\phi\in\mathcal{D}(I)$. In other words, $(u,\psi)=0$ for every test function $\psi$ which is the derivative of a test function. It is easy to see that $\psi$ is the derivative of a test function iff $\int_I\psi(x)\,\mathrm{d}x=0$. Let now $\phi_0$ be any test function with unit integral. Then any $\phi\in\mathcal{D}(I)$ can be decomposed as
\begin{equation}\label{eq:test-function-one-dimensional-zero-integral-decomposition}
\phi(x)=\phi_0(x)\int_I\phi(s)\,\mathrm{d}s+\psi(x),
\end{equation}
where the integral of $\psi$ is zero. Consequently,
\begin{equation}\label{eq:zero-derivative-distribution-constant-pairing}
(u,\phi)=(u,\phi_0)\int_I\phi(x)\,\mathrm{d}x,
\end{equation}
hence $u$ is equal to the constant $(u,\phi_0)$.

We next consider the case where $\Omega$ is a product of intervals: $\Omega=(a_1,b_1)\times(a_2,b_2)\times\cdots\times(a_m,b_m)$. In this case, let $\phi_i\in\mathcal{D}(a_i,b_i)$ be a one-dimensional test function with unit integral. An arbitrary $\phi\in\mathcal{D}(\Omega)$ is now decomposed as follows:
\begin{equation}\label{eq:test-function-product-domain-first-decomposition}
\phi(x_1,\ldots,x_m)=\phi_1(x_1)\int_{a_1}^{b_1}\phi(s_1,x_2,\ldots,x_m)\,\mathrm{d}s_1+\psi_1(x_1,\ldots,x_m).
\end{equation}
The function $\psi_1$ now has the property that
\begin{equation}\label{eq:test-function-first-remainder-zero-integral}
\int_{a_1}^{b_1}\psi_1(x_1,x_2,\ldots,x_m)\,\mathrm{d}x_1=0
\end{equation}
for every $(x_2,\ldots,x_m)$; hence
\begin{equation}\label{eq:test-function-first-remainder-primitive}
\int_{a_1}^{x_1}\psi_1(s,x_2,\ldots,x_m)\,\mathrm{d}s
\end{equation}
is again a test function. Since $\partial u/\partial x_1=0$, it follows that $(u,\psi_1)=0$. Next, we write
\begin{equation}\label{eq:test-function-product-domain-second-decomposition}
\begin{aligned}
&\phi_1(x_1)\int_{a_1}^{b_1}\phi(s_1,x_2,\ldots,x_m)\,\mathrm{d}s_1\\
&\quad=\phi_1(x_1)\phi_2(x_2)\int_{a_1}^{b_1}\int_{a_2}^{b_2}\phi(s_1,s_2,x_3,\ldots,x_m)\,\mathrm{d}s_2\,\mathrm{d}s_1\\
&\qquad+\phi_1(x_1)\psi_2(x_2,\ldots,x_m),
\end{aligned}
\end{equation}
where now
\begin{equation}\label{eq:test-function-second-remainder-zero-integral}
\int_{a_2}^{b_2}\psi_2(x_2,\ldots,x_m)\,\mathrm{d}x_2=0,
\end{equation}
and hence $(u,\phi_1\psi_2)=0$. Proceeding thusly, we finally obtain
\begin{equation}\label{eq:zero-gradient-distribution-product-pairing}
(u,\phi)=(u,\phi_1\phi_2\cdots\phi_m)\int_\Omega\phi(\bm{x})\,\mathrm{d}\bm{x},
\end{equation}
i.e., $u$ is a constant.

For general $\Omega$, it follows from the result just proved that every point has a neighborhood in which $u$ is constant, and of course the constants must be the same if two neighborhoods overlap (Problem~\ref{ex:regular-distributions-injective}). The rest follows from Problem~\ref{ex:distribution-local-vanishing}.
\end{Proof}
